\documentclass[11pt]{article}
\usepackage[margin=1in]{geometry}
\usepackage[T1]{fontenc}
\usepackage{lmodern,amsmath,amssymb,amsthm,mathtools,booktabs,array,longtable,graphicx}
\usepackage{microtype,url,hyperref}
\hypersetup{hidelinks,pdftitle={Contextuality and Exact Information Resources over Finite Chain Rings}}
\newtheorem{theorem}{Theorem}[section]
\newtheorem{lemma}[theorem]{Lemma}
\newtheorem{proposition}[theorem]{Proposition}
\newtheorem{corollary}[theorem]{Corollary}
\theoremstyle{definition}\newtheorem{definition}[theorem]{Definition}
\theoremstyle{remark}\newtheorem{remark}[theorem]{Remark}
\newcommand{\F}{\mathbb F}
\newcommand{\A}{\F_2[u]/(u^4)}
\DeclareMathOperator{\GL}{GL}
\DeclareMathOperator{\Mat}{Mat}
\DeclareMathOperator{\rank}{rank}
\DeclareMathOperator{\diag}{diag}
\DeclareMathOperator{\Span}{span}
\setlength{\emergencystretch}{3em}
\title{Contextuality and Exact Information Resources\\over Finite Chain Rings}
\author{Brian Theory\thanks{Research draft. Author metadata and venue remain to be finalized.}}
\date{Research draft, 19 September 2026}
\begin{document}
\maketitle
\begin{abstract}
We study nonzero bipartite matrices over a finite commutative chain ring using all primitive projective local bases and nonzero-amplitude support. Two Smith invariants determine the support type: the number $r$ of nonzero factors and the multiplicity $h$ of the least exponent. The model is local when $r=1$, logically but not strongly contextual when $r\geq2$ and $h=1$, and strongly contextual when $h\geq2$. A residue-hyperplane criterion and a valuation-uniform Hardy construction give the classification. For a square $d$-dimensional resource, invertible local orbit encoding and one arbitrary complete joint basis readout admit exactly $dh$ messages. Universal exact raw-vector teleportation, with complete rank-one analysis and unrestricted invertible corrections, instead requires an invertible resource. We prove the existence of complete invertible-matrix analyzers over finite local rings and exhibit nonprimitive resources separating maximal orbit coding from teleportation. Tensor choice, preparation closure and a conditional symbolic Boolean interface delimit the interpretation. Finite checks accompany the proofs; no physical probability model or priority claim follows from them.
\end{abstract}

\section{Introduction}
Replacing complex amplitudes by a finite scalar system makes it possible to separate algebraic support from probability. In modal quantum theory, an effect is possible when its contraction with a state is nonzero. Field models already support entanglement, nonlocal support and communication protocols \cite{Schumacher2010,Schumacher2012}. Passing from a field to a ring introduces a different question: which aspects of a resource survive when nonzero amplitudes can annihilate one another?

This paper answers that question for static bipartite resources over finite commutative chain rings, with complete local basis families. The classification uses two elementary pieces of Smith data. Nonzero Smith rank detects whether a Hardy obstruction is available. Leading multiplicity detects whether any global outcome assignment survives. Under a separately stated exact communication task, the same leading multiplicity measures the number of linearly distinguishable orbit messages. Universal exact teleportation distinguishes yet another property, invertibility.

The contribution is the classification and the comparison under matched contracts. Smith theory, the language of support contextuality, Hardy implications and modal protocols are established ingredients. The results do not propose a physical theory, and do not infer a Born rule. In particular, allowing all nonzero ring vectors is useful for static resource questions but does not define a preparation theory closed under arbitrary balanced tensor products.

The proofs are self-contained. The accompanying priority ledger distinguishes a proof from a claim to precedence: a bounded primary-source review found close ingredients and different-context results, but does not certify that no equivalent theorem exists. The scalar-extension appendix connects the results to Boolean correspondence matrices without making either companion manuscript a dependency of the main theorems.

\subsection{Relation to established work}
Schumacher and Westmoreland's two-mobit model already gives a nonlocal support table and a four-message orbit code \cite{Schumacher2010}; their later treatment develops generalized field effects and modal probability-completion questions \cite{Schumacher2012}. These supply background, not new claims here. The support hierarchy follows the global-section formulation of Abramsky and Brandenburger \cite{AbramskyBrandenburger2011}. Our settings are local bases. We do not additionally require the same effect ray to carry a context-independent truth value across distinct bases.

That distinction matters for comparisons with single-system modal Kochen--Specker arguments \cite{SchumacherFreeWill} and the uncolorability of integer-matrix idempotents \cite{BenZvi2017}. The latter concerns compatible idempotent decompositions and transports to arbitrary scalar rings; it is not a Smith classification of a fixed bipartite state with basis-indexed outcomes. Finite-chain-ring module structure and projective ring geometry are likewise classical \cite{Byrne2021,Honold2009}. The use of primitive rows is a standard local-ring construction. Our assertions concern their complete bilinear support tables and two specified information tasks.

State admissibility is a substantive difference from other ring models. In de Beaudrap's generic state space, the coordinates must generate the unit ideal \cite{deBeaudrap2014}. Over a local ring this means that at least one coordinate is a unit. Our nonprimitive nilpotent separator is therefore excluded from that state space; its conclusions belong to the static all-nonzero-vector contract used here. Categorical finite-field protocols can also use a different measurement notion: Varughese's Example~8.3 uses three Bell outcomes on a nine-dimensional pair \cite{Varughese2009}. It does not provide the complete nine-row rank-one analyzer required by our teleportation contract.

\section{Scalar, state and measurement contracts}\label{sec:contract}
Let $K$ be a finite commutative chain ring with identity, maximal ideal $J=(\pi)$, residue field $k=K/J$ of size $q$, and nilpotency index $s$, so $\pi^s=0$ and $\pi^{s-1}\ne0$. The field case has $s=1$. Bars denote reduction modulo $J$. An element is a unit exactly when its residue is nonzero. A square matrix is invertible exactly when its residue is invertible, by the determinant criterion.

Every nonzero element is a unit times a unique power $\pi^a$, $0\leq a<s$. Repeatedly moving an entry of minimal valuation to the first position and using elementary row and column operations gives Smith form
\[
 PMQ^T=\diag(\pi^{a_1},\ldots,\pi^{a_r},0,\ldots),\qquad
 0\leq a_1\leq\cdots\leq a_r<s,
\]
with $P\in\GL_m(K)$ and $Q\in\GL_n(K)$. Clearing a pivot is possible because it divides all remaining entries; induction gives the diagonal form. The invariant exponents can equivalently be recovered from the module type \cite{Byrne2021}. We write $r$ for the number of nonzero factors and
\[
 a=a_1,\qquad h=\#\{i:a_i=a\}.
\]
Writing $M=\pi^a N$ involves a choice of preimage, not division by a unit. Its residue $\bar N$ is well-defined: the ambiguity lies in $(\pi^{s-a})$, which is contained in $J$. Thus $h=\rank_k\bar N$.

\begin{definition}[Complete-basis support]
A primitive row in $K^m$ has a unit coordinate. Its projective ray is its orbit under unit scaling. A local setting is an unordered basis of primitive projective rows, represented by an invertible matrix. For a nonzero resource $M\in K^m\otimes_K K^n\cong\Mat_{m\times n}(K)$, the outcome pair represented by rows $x,y$ is supported exactly when
\[
 xMy^T\ne0.
\]
All such local settings are admitted, and local reversible changes are the full groups $\GL_m(K)$ and $\GL_n(K)$.
\end{definition}

Unit rescaling does not change support. Each context has a supported event because invertible changes cannot send nonzero $M$ to zero. The support marginal for $x$ is independent of Bob's basis: some coordinate of $xMB^T$ is nonzero if and only if $xM\ne0$. Hence these support tables satisfy possibilistic no-signalling, without any assigned real probabilities.

A \emph{global assignment} chooses one outcome for each local setting, with every chosen cross pair supported. A model is \emph{local} here if every supported event extends to such an assignment. It is \emph{logically contextual} if some event does not extend, and \emph{strongly contextual} if no global assignment exists. This is a support-level definition; ``local'' is not a statement about a particular probability distribution.

Two further tasks use the same scalars but different measurements. Their distinctions are essential.
\begin{center}\small
\begin{tabular}{p{.15\linewidth}p{.74\linewidth}}
\toprule
Task & Contract \\\midrule
Support & Nonzero static $M$; balanced tensor over $K$; all primitive projective local bases; nonzero contraction means possible. \\
Orbit coding & Square $M$, $d\geq2$; encodings $UM$ with $U\in\GL_d(K)$; one invertible $K$-linear joint decoder on $d^2$ coordinates; every possible outcome identifies the message. No discard and reprepare. \\
Teleportation & Every nonzero raw input, including correlations with an untouched reference; complete joint rank-one effect basis; all possible branches corrected exactly by allowed invertible maps. No selected successful branch or postselection. \\\bottomrule
\end{tabular}
\end{center}
For support and coding, multiplying a whole resource by a unit is immaterial. Teleportation below demands actual vector recovery and does not silently replace this by an unspecified projective task.

\section{Complete-basis contextuality}\label{sec:classification}
The first lemma requires only that $K$ be finite, commutative and local.
\begin{lemma}[Residue-hyperplane criterion]\label{lem:hyperplane}
A complete-basis support model has a global assignment if and only if there are hyperplanes $H_A\subset k^m$ and $H_B\subset k^n$ such that $xMy^T\ne0$ for every pair of primitive lifts whose residues lie outside $H_A$ and $H_B$.
\end{lemma}
\begin{proof}
Let $U_A$ be the union of the rays selected by a global assignment. Its complement contains no basis, since the corresponding setting would have to select a ray in that complement. If the reductions of the complement spanned $k^m$, one could select a residue basis from them; its lifts would form an invertible basis over $K$. Thus those reductions lie in a proper subspace, which is contained in a hyperplane $H_A$. Every primitive lift of a residue vector outside $H_A$ belongs to $U_A$. Construct $H_B$ in the same way. Every pair from $U_A\times U_B$ is selected in some pair of settings, and is therefore supported.

Conversely, any basis has a row reducing outside a given hyperplane. Select one in every setting on each side. The stated condition makes every chosen cross pair supported, giving a global assignment.
\end{proof}

\begin{lemma}[Uniform Hardy witness]\label{lem:hardy}
For $D=\diag(p,q_0)$ with $p=\pi^a$, $q_0=\pi^b\ne0$ and $a\leq b$, put $c=\pi^{b-a}$. The local bases
\[
 F=\begin{pmatrix}0&1\\1&0\end{pmatrix},\quad
 X=\begin{pmatrix}1&0\\1&1\end{pmatrix},\quad
 C=\begin{pmatrix}1&0\\-c&1\end{pmatrix}
\]
give a supported event with no global extension, using $F,C$ for Alice and $F,X$ for Bob.
\end{lemma}
\begin{proof}
All three determinants are units, and $pc=q_0$. Direct multiplication gives
\begin{align*}
 FDF^T&=\begin{pmatrix}q_0&0\\0&p\end{pmatrix},&
 FDX^T&=\begin{pmatrix}0&q_0\\p&p\end{pmatrix},\\
 CDX^T&=\begin{pmatrix}p&p\\-q_0&0\end{pmatrix},&
 CDF^T&=\begin{pmatrix}0&p\\q_0&-q_0\end{pmatrix}.
\end{align*}
Index outcomes by $0,1$. The supported seed $(F_A=0,F_B=0)$ forces $X_B=1$, then $C_A=0$, and then $F_B=1$, a contradiction. The support strings, in row order, are $1001,0111,1110,0111$. The minus sign cannot be dropped outside characteristic two.
\end{proof}

\begin{theorem}[Smith trichotomy]\label{thm:trichotomy}
For nonzero $M\in\Mat_{m\times n}(K)$, the complete-basis support type is
\[
\begin{array}{rcl}
r=1&\Longleftrightarrow&\text{local},\\
r\geq2,\ h=1&\Longleftrightarrow&\text{logical but not strong},\\
h\geq2&\Longleftrightarrow&\text{strong}.
\end{array}
\]
\end{theorem}
\begin{proof}
Invertible left and right changes merely relabel the complete basis families, so take Smith form.

Suppose $r=1$, with sole nonzero entry $p$ in position $(1,1)$. Retain the outcomes of any supported seed event $(x,y)$. In every other setting choose a row with unit first coordinate. A new--new pair has contraction equal to $p$ times a unit. The seed condition $px_1y_1\ne0$ implies $px_1\ne0$ and $py_1\ne0$, so a seed--new pair is also supported. The seed therefore extends.

If $h=1$, choose a row with unit first coordinate in every setting. Factoring out $\pi^a$, each selected contraction is a product of units plus terms in the radical, hence is a unit. Multiplication by the nonzero $\pi^a$ gives a supported cross pair. A global assignment exists.

If $h\geq2$, write $M=\pi^a M_0$, so $\bar M_0$ has rank $h$. Suppose the hyperplanes of Lemma~\ref{lem:hyperplane} existed, and write $H_B=\ker\beta$. The vectors outside $H_A$ span $k^m$. Since the row space of $\bar M_0$ has dimension at least two, some $x\notin H_A$ has $w=x\bar M_0$ outside the line spanned by $\beta$. There is $y\in\ker w$ with $\beta(y)\ne0$. Lift $x,y$ to rows $\tilde x,\tilde y$. The row $\tilde xM_0$ has a unit coordinate, whereas its contraction with $\tilde y$ lies in $J$. Subtracting that contraction divided by the unit from the corresponding coordinate of $\tilde y$ cancels it exactly and changes no residue. This yields lifts outside both hyperplanes with zero contraction, a contradiction.

Finally, if $r\geq2$, embed Lemma~\ref{lem:hardy} in the first two nonzero Smith coordinates and extend its bases with standard rows. At each forcing step, the forced row has zero contraction with every added coordinate outcome. Thus the same seed cannot extend, even in higher dimensions. These observations exhaust all nonzero cases.
\end{proof}

There are $q^{sm}-q^{(s-1)m}$ primitive rows in $K^m$. Unit scaling acts freely, and $|K^\times|=q^{s-1}(q-1)$. Consequently the numbers of rays and unordered projective bases are
\[
 q^{(s-1)(m-1)}\frac{q^m-1}{q-1},\qquad
 \frac{q^{m^2(s-1)}|\GL_m(k)|}{|K^\times|^m m!}.
\]
For $K=\A$ and $m=2$, these are $24$ rays and $192$ bases. These counting facts are background; enumeration is unnecessary for the classification proof.

\section{Exact orbit coding and teleportation}\label{sec:tasks}
\begin{lemma}[Invertible matrices span]\label{lem:span}
For a field $k$ and $d\geq2$, $\GL_d(k)$ spans $\Mat_d(k)$. Over a finite commutative local ring $K$, there is a $K$-basis of $\Mat_d(K)$ consisting entirely of invertible matrices.
\end{lemma}
\begin{proof}
For $i\ne j$, $E_{ij}=(I+E_{ij})-I$. For a diagonal unit $E_{ii}$, choose a permutation matrix $P$ with $Pe_j=e_i$, $j\ne i$. Then $P(I+E_{ji})-P=E_{ii}$, again a difference of invertibles. This works over every field, including $\F_2$. Select $d^2$ independent invertibles from this spanning family over $k$ and lift their entries to $K$. Each lifted matrix is invertible. The square matrix of their flattened vectors is also invertible, because its residue is a basis matrix. The lifts therefore form the asserted basis.
\end{proof}

\begin{theorem}[Leading-layer orbit codebook]\label{thm:code}
For $M\in\Mat_d(K)\setminus\{0\}$, $d\geq2$, the maximum number of messages under the orbit-coding contract of Section~\ref{sec:contract} is $N_{\max}=dh$.
\end{theorem}
\begin{proof}
Write $M=\pi^aN$. The leading residue of each divided codeword $UN$ lies in
\[
 V=\{X\bar N:X\in\Mat_d(k)\},\qquad \dim_k V=dh.
\]
Every orbit codeword has a nonzero leading residue. An invertible joint decoder preserves this fact. Exact decoding forces supports for different messages to be disjoint; their nonzero leading residue vectors are therefore independent. Hence their number is at most $dh$.

By Lemma~\ref{lem:span}, images of invertible residue matrices span $V$. Choose $dh$ of these images that are independent and lift the encoders to $U_1,\ldots,U_{dh}\in\GL_d(K)$. The columns $\operatorname{vec}(U_iN)$ are independent modulo $J$. Extend the residue columns to a basis and lift the added columns arbitrarily. The resulting matrix $B$ over $K$ is invertible. Its inverse sends the actual codewords to $\pi^ae_i$, each with a unique possible coordinate outcome. This attains the bound.
\end{proof}

\begin{corollary}\label{cor:equiv}
With the complete-basis support and orbit-coding contracts matched, $N_{\max}>d$ if and only if $M$ is strongly contextual.
\end{corollary}
\begin{proof}
Both conditions are equivalent to $h\geq2$.
\end{proof}
The $d$-message comparison is the dimension baseline for a fresh $d$-coordinate carrier with complete basis readout. It is not an entropy or an asymptotic channel capacity. If encoders may discard the resource, or decoders are restricted, the stated optimization changes. Over a field $h=r$, giving the familiar linear-algebra orbit count $dr$; the two-mobit four-message protocol is already known \cite{Schumacher2010}.

\begin{theorem}[Universal raw teleportation]\label{thm:teleport}
Let $K$ be a finite commutative local ring. For nonzero $M\in\Mat_d(K)$, $d\geq2$, universal exact raw-vector teleportation with complete rank-one analysis and unrestricted invertible corrections is possible if and only if $M$ is invertible.
\end{theorem}
\begin{proof}
Fix vectorization so that the branch labeled by the reshaped joint effect $E$ transfers the input by $T=M^TE^T$. If a nonzero branch map had a nonzero kernel vector $z$, choose $x$ with $Tx\ne0$. Both $x$ and $x+z$ are nonzero and give the same possible branch output. A fixed correction cannot recover both. Thus each nonzero universally correct branch is injective. On the finite free module $K^d$, it is bijective and its inverse is linear; its invertibility forces $M$ to be invertible. A complete effect basis and nonzero $M$ exclude the possibility that every branch map is zero: the reshapes span all matrices, so this would force $M^TE_{ij}=0$ for all $i,j$ and hence $M=0$.

Conversely, choose the invertible-matrix basis of Lemma~\ref{lem:span} as the reshapes of the joint effects. It is a complete analysis basis, and every transfer $M^TE^T$ is invertible. Correct with its inverse. The identity $T^{-1}T=I$ remains an identity after tensoring with the identity on an untouched reference, proving the reference-preserving statement as well.
\end{proof}

The analyzer existence step extends the explicitly constructed dimension-two and dimension-eight analyzers in the supplied research. It is a residue-lifting corollary of elementary matrix algebra. Its inclusion removes a dimension restriction from sufficiency when all the required bases and corrections are admitted; it does not assert that a physically restricted gate set can implement them.

\begin{corollary}[A nonprimitive separator]\label{cor:separator}
For a chain ring of length $s>1$ and $0<a<s$, the resource $\pi^a I_d$ is strongly contextual and has $d^2$ exact orbit messages, but cannot teleport universally under Theorem~\ref{thm:teleport}.
\end{corollary}
\begin{proof}
Its leading multiplicity is $d$. It also annihilates every vector multiplied by $\pi^{s-a}$, so is singular. Apply Theorems~\ref{thm:trichotomy}--\ref{thm:teleport}.
\end{proof}
The smallest example is $\varepsilon I_2$ over $\F_2[\varepsilon]/(\varepsilon^2)$. This resource is nonprimitive. Restricting preparation to primitive vectors excludes this particular separator and must be stated explicitly.

\section{Boundaries of the resource comparison}\label{sec:boundaries}
\subsection{Balanced tensors and regular binary expansion}
In $A=\A$, balanced scalar multiplication identifies $A\otimes_A A$ with $A$. Therefore
\[
 u^3\otimes_Au=0,\qquad u^3\otimes_{\F_2}u\ne0.
\]
The second assertion follows by applying field-linear functionals that are nonzero on each factor. Nonzero preparations thus need not remain nonzero under the balanced tensor. Primitive vectors do remain primitive, since a product of unit coordinates is a unit. Yet that restriction is not closed under the proposed conditional support update: applying the second standard effect to $\diag(1,u)$ gives the nonprimitive nonzero vector $(0,u)$.

The regular representation of $A$ is a faithful algebra map to $4\times4$ binary matrices. It is not a monoidal equivalence between the balanced and literal tensor spaces. Expanding a Smith factor $u^a$ gives binary rank $4-a$, whence
\[
 \rank_{\F_2}\lambda(M)=\sum_{i=1}^r(4-a_i).
\]
Unrestricted independent binary row and column operations classify a matrix only by that total rank. They can consequently erase distinctions detected by restricted $A$-linear operations.

For instance, $\diag(1,u^2)$ and $\diag(u,u)$ both have binary rank six but leading multiplicities one and two. To avoid relying on a nonprimitive comparison, the primitive matrices
\[
 \diag(1,1,u^3),\qquad \diag(1,u,u^2)
\]
both have binary rank nine, while their leading multiplicities are two and one. Their $A$-linear orbit codebooks have sizes six and three; both are singular, and hence neither universally teleports in the stated task.

In one audited restricted bipartite construction, shared and literal support tables agree after transposing Bob's regular blocks and applying the involution $R\mapsto R^{-1}$. This is an isomorphism of those particular outcome tables. It does not identify their tensor spaces or preparation rules. Figure~\ref{fig:tensors} records these distinct comparisons.

\begin{figure}[ht]\centering
\fbox{\begin{minipage}{.91\linewidth}
\textbf{Three comparisons with different contracts}\par\smallskip
Balanced scalar tensor: $u^3\otimes_Au=0$; literal field tensor: $u^3\otimes_{\F_2}u\ne0$.\par\smallskip
Restricted $A$ operations retain Smith data; unrestricted binary operations retain total binary rank.\par\smallskip
An outcome-table relabeling for selected measurements preserves those supports only. It implies neither equality of tensor spaces nor unrestricted operational equivalence.
\end{minipage}}
\caption{Tensor choice, operation group and table equivalence are separate statements.}\label{fig:tensors}
\end{figure}

\subsection{A boundary on the number of settings}
\begin{theorem}\label{thm:two}
A nonzero tensor with any finite number of parties over $\F_2$, two-dimensional local carriers, and at most two complete binary settings per party has a global assignment. The same conclusion holds over a finite commutative local ring with residue field $\F_2$.
\end{theorem}
\begin{proof}
Two distinct binary bases share a row $a$ and have other rows $b$ and $a+b$. If they coincide, use $b$ as the other row in both settings. Expand the tensor in the coordinates dual to $(a,b)$ at each site. Choose an inclusion-minimal set $S$ of sites using the $b$ coordinate whose coefficient is nonzero. Outside $S$ select $a$. Inside $S$ select the nonshared row of the chosen setting. Every joint contraction is the coefficient for $S$ plus coefficients of proper subsets, all of which vanish. Thus every selected joint event is supported.

For a local ring with radical $J$, choose the last nonzero radical layer containing all tensor coefficients: the tensor lies in $(J^a)^{2^n}$ but not in $(J^{a+1})^{2^n}$. Choose an $\F_2$-linear functional on $J^a/J^{a+1}$ for which the coefficient tensor is nonzero. Reduce each local basis modulo $J$ and apply the field argument. Every selected original contraction has nonzero image under this functional and is therefore nonzero.
\end{proof}
This does not exclude Hardy logical contextuality, which concerns failure to extend a particular event. Nor does it cover three settings or a different residue field. For two parties, the exclusion of a strongly contextual binary two-setting table is close to the established PR-box and modal-table boundaries \cite{AbramskyBrandenburger2011,Schumacher2012}; the proof above states explicitly what extends to arbitrary party number and radical layers.

\section{Finite verification and reproducibility}\label{sec:verification}
The analytic proofs determine the general statements. Independent finite checks examine whether the implementations and selected examples agree with those statements. They do not establish mathematical priority or replace quantifiers in a proof.

\begin{table}[ht]\centering\small
\begin{tabular}{rrrrrrll}
\toprule
$(a,b)$ & Count & $r$ & $h$ & Binary rank & Messages & Support & Teleport\\\midrule
$(0,0)$ & 24,576 & 2 & 2 & 8 & 4 & Strong & Yes\\
$(0,1)$ & 18,432 & 2 & 1 & 7 & 2 & Logical & No\\
$(0,2)$ & 9,216 & 2 & 1 & 6 & 2 & Logical & No\\
$(0,3)$ & 4,608 & 2 & 1 & 5 & 2 & Logical & No\\
$(0,4)$ & 4,608 & 1 & 1 & 4 & 2 & Local & No\\
$(1,1)$ & 1,536 & 2 & 2 & 6 & 4 & Strong & No\\
$(1,2)$ & 1,152 & 2 & 1 & 5 & 2 & Logical & No\\
$(1,3)$ & 576 & 2 & 1 & 4 & 2 & Logical & No\\
$(1,4)$ & 576 & 1 & 1 & 3 & 2 & Local & No\\
$(2,2)$ & 96 & 2 & 2 & 4 & 4 & Strong & No\\
$(2,3)$ & 72 & 2 & 1 & 3 & 2 & Logical & No\\
$(2,4)$ & 72 & 1 & 1 & 2 & 2 & Local & No\\
$(3,3)$ & 6 & 2 & 2 & 2 & 4 & Strong & No\\
$(3,4)$ & 9 & 1 & 1 & 1 & 2 & Local & No\\
$(4,4)$ & 1 & 0 & 0 & 0 & n/a & Excluded & n/a\\
\bottomrule\end{tabular}
\caption{All fifteen Smith classes over $A=\mathbb F_2[u]/(u^4)$ for $2\times2$ matrices. Exponent $4$ denotes zero. Logical means logical but not strong. Counts include the zero matrix; the resource tasks exclude it.}\label{tab:smith}
\end{table}


Table~\ref{tab:smith} is regenerated from the supplied, independently rechecked inventory. For $A^2$, the $65{,}535$ nonzero matrices divide into $5{,}265$ local, $34{,}056$ logical-but-not-strong and $26{,}214$ strong instances. The inventory includes zero separately, outside the resource model. A separate direct support/2-SAT implementation over the dual numbers classifies all $255$ nonzero $2\times2$ matrices as $81$ local, $72$ logical-but-not-strong and $102$ strong. That method does not use the hyperplane classification as its decision rule. The supplement also contains explicit encoders and joint decoder matrices for all fourteen nonzero $A$ Smith representatives.

\begin{center}\small
\begin{tabular}{p{.25\linewidth}p{.61\linewidth}}
\toprule
Evidence & Scope and limitation\\\midrule
Analytic proofs & All finite commutative chain rings for classification/coding; all finite commutative local rings for analyzer lifting and raw teleportation.\\
Independent enumeration & $2\times2$ resources over $A$ and the dual numbers; finite domains only.\\
Explicit certificates & Fourteen nonzero $A$ Smith representatives, with actual invertible encoder/decoder matrices.\\
Supplied historical suites & Separately rerun in an isolated directory; assertion counts are not merged with exhaustive-state counts.\\
Symbolic bridge checks & All $256$ products in the fixed sixteen-element LM family; $112$ remain inside and $144$ do not.\\\bottomrule
\end{tabular}
\end{center}

The supplement records the unchanged canonical runner, exact commands, interpreter and dependency versions, input/output SHA-256 hashes, runtime and per-stage status. It also records the new analyzer argument and a finite verification of its two independent invertibility requirements. No new timing campaign is part of these mathematical checks.

\section{Discussion}
The resource hierarchy is controlled by different questions about the same Smith form. Nonzero rank permits a logical obstruction; leading multiplicity controls the absence of every global assignment and the size of an exact orbit code; invertibility controls universal raw-vector transfer. None of these statements identifies a physical capability without the relevant operation and readout assumptions.

There are two particularly useful limits. First, the nilpotent maximal-code example is a static nonprimitive resource, so it cannot by itself justify a closed process theory. Second, the codebook/contextuality equivalence compares matched complete-basis contracts. Restricting gates, changing the scalar tensor, adding discard-and-reprepare, or weakening exactness requires a new theorem.

The next mathematical questions are a classification beyond chain rings, a process semantics closed under both preparation and conditioning, and resource comparisons for restricted measurement families. These remain open here. The publication priority of the specific classification and exact orbit formula also remains subject to the unresolved source checks listed in the accompanying ledger; the paper claims the displayed proofs, not certified historical firstness.

\appendix
\section{A conditional interface from Boolean formulas}\label{app:lm}
Let $B$ be the Boolean algebra of formulas modulo logical equivalence, regarded as a Boolean ring over $\F_2$. Paper B's formula-valued matrix is
\[
 M_f(X,Y)=\begin{pmatrix}f(X,Y)&f(X,\neg Y)\\f(\neg X,Y)&f(\neg X,\neg Y)\end{pmatrix}.
\]
Its valuation and pairing identities belong to that Boolean calculus \cite{PaperB}. They do not choose amplitudes, a tensor base, effects or a possibility rule. In particular, every Boolean-ring element is idempotent, while the only idempotents in a local ring are $0,1$: if $e(1-e)=0$, one of $e,1-e$ is a unit. Thus a unital Boolean-ring map into $A$ has image in $\{0,1\}$ and cannot generate $u$.

Choose instead the scalar extension $B_A=B\otimes_{\F_2}A$. Each element is uniquely $\sum_{j=0}^3 f_j u^j$. A Boolean valuation $v:B\to\F_2$ extends coefficientwise to $v_A:B_A\to A$. For a selected enriched state $\psi$ and effect $e$, write $e\psi=\sum_j g_j u^j$ and define the Boolean formula
\[
 \operatorname{Supp}(e\psi)=g_0\lor g_1\lor g_2\lor g_3.
\]
Then
\[
 v(\operatorname{Supp}(e\psi))=1
 \quad\Longleftrightarrow\quad v_A(e)v_A(\psi)\ne0.
\]
This follows from coefficientwise valuation and uniqueness of the $u$ expansion. The disjunction is essential; nonzero testing is neither an additive nor a multiplicative homomorphism. The construction is natural relative to the chosen enrichment and valuation, not a unique modal interpretation of bare Boolean formulas. It applies to the specified $\F_2$-algebra $A$, not automatically to mixed-characteristic $K$.

For an invertible coordinate readout $B_0$, put $J_i=B_0^{-1}P_iB_0$, with $P_i$ the standard coordinate projections. Direct multiplication yields $J_i^2=J_i$, $J_iJ_j=0$ for $i\ne j$, and $\sum_iJ_i=I$. Moreover $J_i\psi$ is nonzero exactly when the measured coordinate $(B_0\psi)_i$ is nonzero, because the corresponding column of $B_0^{-1}$ has a unit coordinate. No division by that measured amplitude is required. Symbolic valuation commutes with these operations when $B_0$ is invertible symbolically. This supplies repeatable linear branches, but no probability normalization or repair of primitive-state closure.

\begin{figure}[ht]\centering
\fbox{\begin{minipage}{.91\linewidth}
Boolean formulas $B$ and their valuations are the starting data.\par\smallskip
$B\longrightarrow B\otimes_{\F_2}A$: scalar enrichment requires the explicit choice of $A$.\par\smallskip
Choose carrier, tensor base, gates and effects. Coefficientwise valuation then produces $A$-valued contractions.\par\smallskip
Apply coefficient OR to obtain symbolic support. No arrow in this construction supplies real probabilities or a physical realization.
\end{minipage}}
\caption{The typed interface distinguishes proved maps from additional modeling choices.}\label{fig:interface}
\end{figure}

\section{Resource map}
The table below collects the square-resource cases under the contracts of Section~\ref{sec:contract}.
\begin{figure}[ht]\centering
\begin{tabular}{ccl}
\toprule
Smith region & Exact orbit count & Support and teleportation\\\midrule
$r=1$ & $d$ & Local; singular for $d\geq2$\\
$r\geq2,\ h=1$ & $d$ & Logical, not strong; singular\\
$2\leq h<d$ & $dh$ & Strong; singular\\
$h=d,\ a>0$ & $d^2$ & Strong; nonprimitive, singular\\
$h=d,\ a=0$ & $d^2$ & Strong; invertible, universally teleporting\\\bottomrule
\end{tabular}
\caption{Square-resource map under the stated contracts, $d\geq2$. Zero is excluded. The nilpotent separator occupies the fourth row; the primitive rank-nine pair occupies the third and second rows with $d=3$.}\label{fig:map}
\end{figure}
\bibliographystyle{plain}
\bibliography{references}
\end{document}
